\documentclass[../../main.tex]{subfiles}

\begin{document}

\chapter{Analisis Kode Mesin, Disassembly, dan Optimasi}
\label{chap:disassembly}

% =========================================================
% OUTCOME PEMBELAJARAN KHUSUS BAB
% =========================================================

\begin{subcpmkbox}
\SubCPMK{4.3} --- Menganalisis kode mesin hasil \eng{disassembly} dan
menerapkan strategi optimasi program pada level instruksi.

\vspace{0.4em}
\textbf{CPMK yang didukung:} \CPMK{04} \quad\textbullet\quad
\textbf{CPL yang didukung:} \CPL{03}
\end{subcpmkbox}

\begin{tujuanbox}
Setelah mempelajari bab ini, mahasiswa mampu:
\begin{enumerate}
    \item menjelaskan pengertian \eng{disassembly} dan perbedaannya
          dengan \eng{assembly};
    \item memakai perintah \code{U} pada DEBUG serta perkakas
          sejenis untuk mengubah kode mesin kembali menjadi
          \eng{mnemonic};
    \item menyebutkan keterbatasan \eng{disassembly}, khususnya
          sulitnya membedakan kode dan data serta munculnya alamat
          relatif;
    \item mengenali pola hasil \eng{disassembly}: pembandingan dan
          percabangan, perulangan, \eng{prologue}/\eng{epilogue}
          prosedur, inisialisasi segmen, dan pemanggilan interupsi;
    \item menghitung metrik optimasi (jumlah instruksi, perkiraan
          siklus, ukuran byte, dan jumlah akses memori) serta
          menjelaskan bahwa tanpa pengukuran optimasi hanyalah
          dugaan;
    \item menerapkan teknik optimasi level instruksi: memilih
          register daripada memori, mengganti \asm{MUL}/\asm{DIV}
          dengan pergeseran, memakai \asm{LEA}, mengurangi lompatan,
          \eng{loop unrolling}, dan \lstasm|XOR reg, reg|;
    \item menerapkan \eng{strength reduction} dan memilih di antara
          optimasi ukuran atau optimasi kecepatan sesuai batasan
          program;
    \item merancang pengukuran yang adil untuk membandingkan versi
          naif dan versi teroptimasi.
\end{enumerate}
\end{tujuanbox}

\begin{materibox}
\begin{enumerate}
    \item Pengertian \eng{disassembly}: mengembalikan kode mesin
          menjadi \eng{mnemonic}; perintah \code{U} pada DEBUG dan
          perkakas sejenis.
    \item Keterbatasan \eng{disassembly}: kode dan data sulit
          dibedakan serta munculnya alamat relatif dan displacement.
    \item Pola hasil \eng{disassembly}: pembandingan dan percabangan,
          perulangan, \eng{prologue}/\eng{epilogue} prosedur,
          inisialisasi segmen, dan pemanggilan interupsi.
    \item Metrik optimasi: jumlah instruksi, perkiraan jumlah siklus,
          ukuran kode dalam byte, dan jumlah akses memori.
    \item Teknik optimasi level instruksi: register daripada memori,
          pergeseran pengganti \asm{MUL}/\asm{DIV}, \asm{LEA},
          penyusunan ulang alur, \eng{loop unrolling}, dan
          \lstasm|XOR reg, reg|.
    \item \eng{Strength reduction} dan tabel penggantian operasi.
    \item Trade-off optimasi ukuran dan optimasi kecepatan.
    \item Pengukuran: penghitung iterasi, pencatatan waktu melalui
          layanan sistem, dan perbandingan sebelum/sesudah.
    \item Etika dan batas optimasi: keterbacaan, kemudahan
          pemeliharaan, dan larangan optimasi dini.
\end{enumerate}
\end{materibox}

\subfile{section-01}
\subfile{section-02}
\subfile{section-03}

\section{Contoh Program}
\label{sec:contoh-bab16}

Bagian ini menyajikan tiga program yang mengerjakan pekerjaan
\emph{sama persis}, yaitu menjumlahkan sepuluh elemen larik lalu
mengalikan jumlahnya dengan delapan. Larik berisi
$10, 20, \dots, 100$, sehingga jumlahnya $550$ dan hasil akhirnya
$4400$ atau \hex{1130}. Ketiga program dibedakan hanya oleh cara
mengerjakannya, sehingga keduanya dapat dibandingkan secara adil.

\subsection{Contoh 1: Versi Naif}

Kode~\ref{lst:program-naif} adalah versi yang ditulis tanpa
memikirkan efisiensi. Program ini benar dan mudah dibaca, tetapi
memuat beberapa pemborosan.

\lstinputlisting[
    language=MASM,
    caption={Versi naif: \file{naif.asm}.},
    label={lst:program-naif}
]{\CodeDir/naif.asm}

Empat inefisiensi dapat ditemukan pada program tersebut.

\begin{enumerate}
    \item Instruksi \lstasm|ADD AX, larik[SI]| memakai pengalamatan
          langsung berindeks, sehingga alamat efektifnya memerlukan
          8 siklus, bukan 5 siklus seperti \lstasm|[SI]|.
    \item \lstasm|MOV AX, 0| berukuran tiga byte, padahal ada cara
          yang hanya dua byte.
    \item Hasil akhir ditulis ke memori (\code{jumlah}) lalu dibaca
          kembali, padahal nilainya sudah ada di \reg{AX}. Ini
          menambah dua akses memori yang tidak perlu.
    \item \lstasm|MUL BX| untuk mengalikan dengan delapan adalah
          pilihan paling mahal di seluruh program: 118--133 siklus
          untuk pekerjaan yang dapat diselesaikan dalam enam siklus.
\end{enumerate}

\subsection{Contoh 2: Versi Teroptimasi}

Kode~\ref{lst:program-teroptimasi} memperbaiki keempat hal tersebut
tanpa mengubah hasil akhir.

\lstinputlisting[
    language=MASM,
    caption={Versi teroptimasi: \file{teroptimasi.asm}.},
    label={lst:program-teroptimasi}
]{\CodeDir/teroptimasi.asm}

Empat perbaikan yang diterapkan adalah sebagai berikut.

\begin{itemize}
    \item \lstasm|MOV SI, OFFSET larik| membuat \reg{SI} menunjuk
          langsung ke data, sehingga pengalamatan
          \lstasm|ADD AX, [SI]| hanya memerlukan 5 siklus alamat
          efektif.
    \item \lstasm|XOR AX, AX| mengosongkan akumulator dalam dua byte
          dan tiga siklus.
    \item Nilai antara tidak lagi bolak-balik ke memori; \reg{AX}
          langsung dipakai sebagai akumulator sampai akhir.
    \item \lstasm|MUL BX| diganti tiga kali \lstasm|SHL AX, 1|.
          Totalnya enam siklus, yaitu sekitar dua puluh kali lebih
          cepat daripada satu instruksi \asm{MUL}.
\end{itemize}

\subsection{Contoh 3: Versi Teroptimasi dengan \eng{Loop Unrolling}}

Kode~\ref{lst:program-unroll} menambahkan satu teknik lagi:
isi perulangan digandakan sehingga satu iterasi mengerjakan dua
elemen. Jumlah iterasi turun dari sepuluh menjadi lima.

\lstinputlisting[
    language=MASM,
    caption={Versi teroptimasi dengan \eng{unrolling}:
    \file{teroptimasi-unroll.asm}.},
    label={lst:program-unroll}
]{\CodeDir/teroptimasi-unroll.asm}

\subsection{Perbandingan Hasil Disassembly}

Gambar~\ref{fig:disasm-bandingan} menampilkan potongan hasil
\eng{disassembly} kedua versi utama secara berdampingan. Keduanya
dibongkar dengan asumsi program dimuat pada alamat
\hex{1D2F:0100}, larik berada pada offset \hex{01F4}, variabel
\code{jumlah} pada offset \hex{0208}, dan variabel \code{hasil} pada
offset \hex{020A}.

\begin{figure}[H]
    \centering
    \begin{minipage}[t]{0.49\textwidth}
        \centering
        \textbf{Versi naif}\\[0.3em]
        \begin{lstlisting}[language=MASM, numbers=none, aboveskip=0pt]
1D2F:0100 B90A00   MOV  CX,000A
1D2F:0103 BE0000   MOV  SI,0000
1D2F:0106 B80000   MOV  AX,0000
1D2F:0109 0384F401 ADD  AX,[SI+01F4]
1D2F:010D 83C602   ADD  SI,+02
1D2F:0110 E2F7     LOOP 0109
1D2F:0112 89060802 MOV  [0208],AX
1D2F:0116 8B060802 MOV  AX,[0208]
1D2F:011A BB0800   MOV  BX,0008
1D2F:011D F7E3     MUL  BX
1D2F:011F 89060A02 MOV  [020A],AX
        \end{lstlisting}
        {\scriptsize 11 instruksi, 35 byte}
    \end{minipage}
    \hfill
    \begin{minipage}[t]{0.49\textwidth}
        \centering
        \textbf{Versi teroptimasi}\\[0.3em]
        \begin{lstlisting}[language=MASM, numbers=none, aboveskip=0pt]
1D2F:0100 B90A00   MOV  CX,000A
1D2F:0103 BEF401   MOV  SI,01F4
1D2F:0106 31C0     XOR  AX,AX
1D2F:0108 0304     ADD  AX,[SI]
1D2F:010A 83C602   ADD  SI,+02
1D2F:010D E2F9     LOOP 0108
1D2F:010F D1E0     SHL  AX,1
1D2F:0111 D1E0     SHL  AX,1
1D2F:0113 D1E0     SHL  AX,1
1D2F:0115 89060A02 MOV  [020A],AX
        \end{lstlisting}
        {\scriptsize 10 instruksi, 25 byte}
    \end{minipage}
    \caption{Perbandingan hasil \eng{disassembly} versi naif dan
    versi teroptimasi.}
    \label{fig:disasm-bandingan}
\end{figure}

Tiga perbedaan penting tampak pada kedua kolom tersebut.

\begin{itemize}
    \item \lstasm|ADD AX, [SI+01F4]| pada versi naif mengandung
          \eng{displacement} \hex{01F4}: nama simbolik \code{larik}
          hilang dan berubah menjadi angka. Versi teroptimasi
          menghasilkan \lstasm|ADD AX, [SI]| yang lebih pendek dua
          byte.
    \item \lstasm|MOV AX, 0000| (3 byte) digantikan
          \lstasm|XOR AX, AX| (2 byte).
    \item Urutan \lstasm|MOV [0208],AX|, \lstasm|MOV AX,[0208]|,
          \lstasm|MOV BX,0008|, \lstasm|MUL BX| digantikan oleh tiga
          \lstasm|SHL AX, 1|. Perhatikan bahwa hasil
          \eng{disassembly} sama sekali tidak memberi tahu mengapa
          perubahan itu dilakukan; penjelasannya hanya ada pada
          perbandingan metrik.
\end{itemize}

\subsection{Perbandingan Hasil Pengukuran}

Tabel~\ref{tab:hasil-pengukuran-16} merangkum perbandingan ketiga
versi. Angka siklus dihitung memakai perkiraan pada
Tabel~\ref{tab:siklus-8086-perkiraan} dengan asumsi tanpa
\eng{wait state} dan \eng{prefetch queue} selalu terisi. Bagian yang
diukur adalah bagian yang ditandai komentar pada ketiga berkas.

\begin{table}[H]
    \centering
    \caption{Perbandingan hasil pengukuran versi naif dan versi
    teroptimasi.}
    \label{tab:hasil-pengukuran-16}
    \small
    \begin{tabular}{@{}p{3.6cm}cccc@{}}
        \toprule
        \textbf{Versi} & \textbf{Instruksi} & \textbf{Ukuran} &
        \textbf{Perkiraan} & \textbf{Akses} \\
        & \textbf{dieksekusi} & \textbf{kode} &
        \textbf{siklus} & \textbf{memori} \\
        \midrule
        Naif & 38 & 35 byte & $\approx$546 & 13 \\
        \addlinespace
        Teroptimasi & 37 & 25 byte & $\approx$370 & 11 \\
        \addlinespace
        Teroptimasi + \eng{unrolling} & 27 & 28 byte & $\approx$280
        & 11 \\
        \bottomrule
    \end{tabular}
\end{table}

Tabel~\ref{tab:hasil-pengukuran-16} menyimpan pelajaran penting.
Pertama, jumlah instruksi yang dieksekusi hampir tidak berubah
antara versi naif dan versi teroptimasi (38 berbanding 37), padahal
perkiraan siklusnya turun sekitar 32 persen. Artinya, \important{jumlah
instruksi saja adalah metrik yang menyesatkan}; yang menentukan
adalah \emph{jenis} instruksi dan banyaknya akses memori. Kedua,
versi \eng{unrolling} paling cepat tetapi ukurannya bertambah tiga
byte dibandingkan versi teroptimasi; inilah trade-off yang dibahas
pada Subbab~\ref{sec:optimasi-instruksi}. Ketiga, angka-angka ini
baru perkiraan. Pembuktiannya tetap memerlukan pengukuran waktu pada
mesin, seperti yang akan dikerjakan pada praktikum.

\section{Praktikum}
\label{sec:praktikum-bab16}

\begin{praktikumbox}
\textbf{Judul:} Menganalisis, mengoptimasi, dan mengukur program pada
level instruksi.\\[0.3em]
\textbf{Sub-CPMK:} 4.3 \quad\textbullet\quad
\textbf{Alokasi waktu:} 150 menit \quad\textbullet\quad
\textbf{Bentuk:} kelompok (2--3 mahasiswa) dan laporan individu
\end{praktikumbox}

\subsection{Tujuan Praktikum}

Setelah menyelesaikan praktikum ini, mahasiswa mampu membongkar
program menjadi kode mesin, menemukan inefisiensi berdasarkan metrik,
menulis versi perbaikan, serta membuktikan perbaikannya melalui
pengukuran yang adil.

\subsection{Alat dan Bahan}

\begin{itemize}
    \item Komputer dengan DOSBox dan MS-DEBUG, atau Emu8086.
    \item TASM dan TLINK (atau GUI Turbo Assembler).
    \item Berkas \file{naif.asm}, \file{teroptimasi.asm}, dan
          \file{teroptimasi-unroll.asm} dari \file{code/asm/}.
    \item Lembar kerja pengukuran dan formulir jurnal praktikum.
\end{itemize}

\subsection{Langkah Kerja}

\begin{enumerate}
    \item \textbf{Merakit dan menjalankan tiga versi program.}
    \begin{enumerate}
        \item Rakit ketiga berkas dengan perintah \file{tasm} dan
              \file{tlink} di dalam DOSBox.
        \item Jalankan tiap versi, lalu periksa isi variabel
              \code{hasil} melalui jendela memori Emu8086 atau perintah
              \code{D} pada DEBUG untuk memastikan ketiganya
              menghasilkan \hex{1130}.
        \item Catat panjang berkas \file{.exe} masing-masing versi.
    \end{enumerate}

    \item \textbf{Membongkar kode mesin.}
    \begin{enumerate}
        \item Muat program dengan perintah \code{N} dan \code{L} pada
              DEBUG, lalu bongkar mulai titik masuk program dengan
              perintah \code{U}.
        \item Catat alamat, byte mesin, dan \eng{mnemonic} setiap
              instruksi pada lembar kerja.
        \item Tandai setiap pola yang telah dipelajari pada
              Tabel~\ref{tab:pola-disassembly}: perulangan,
              inisialisasi segmen, dan sebagainya.
        \item Ulangi untuk versi teroptimasi, lalu susun kedua hasil
              secara berdampingan seperti pada
              Gambar~\ref{fig:disasm-bandingan}.
    \end{enumerate}

    \item \textbf{Mengidentifikasi inefisiensi.}
    \begin{enumerate}
        \item Tandai instruksi pada versi naif yang memakai akses
              memori berulang, memuat konstanta panjang, atau
              memakai \asm{MUL}/\asm{DIV}.
        \item Hitung perkiraan siklus bagian yang diukur dengan
              bantuan Tabel~\ref{tab:siklus-8086-perkiraan}.
        \item Tuliskan usulan perbaikan beserta alasan dan syarat
              berlakunya.
    \end{enumerate}

    \item \textbf{Menulis versi perbaikan sendiri.}
    \begin{enumerate}
        \item Salin \file{naif.asm} menjadi \file{optimasi-anda.asm}.
        \item Terapkan minimal tiga teknik dari
              Tabel~\ref{tab:strength-reduction}.
        \item Pastikan keluaran tetap sama: variabel \code{hasil}
              harus tetap bernilai \hex{1130}.
    \end{enumerate}

    \item \textbf{Mengukur secara adil.}
    \begin{enumerate}
        \item Tambahkan penghitung iterasi dan pencatat waktu
              \code{INT 1Ah} fungsi \hex{00} seperti pada
              Subbab~\ref{sec:optimasi-instruksi}.
        \item Ukur ketiga versi (naif, teroptimasi, dan
              \file{optimasi-anda.asm}) dengan larik dan jumlah
              iterasi yang sama.
        \item Ulangi setiap pengukuran lima kali, buang percobaan
              pertama, lalu hitung nilai tengahnya.
        \item Isi Tabel~\ref{tab:hasil-pengukuran-16} dengan hasil
              pengukuran Anda, bukan dengan angka perkiraan.
    \end{enumerate}
\end{enumerate}

\subsection{Hasil yang Dikumpulkan}

\begin{enumerate}
    \item Berkas \file{optimasi-anda.asm} yang dapat dirakit tanpa
          kesalahan.
    \item Tabel perbandingan hasil \eng{disassembly} versi naif dan
          versi perbaikan.
    \item Tabel hasil pengukuran sebelum dan sesudah, lengkap dengan
          perkiraan siklus dan jumlah akses memori.
    \item Analisis singkat: teknik mana yang memberi perbaikan
          terbesar, dan teknik mana yang menambah ukuran kode.
\end{enumerate}

\subsection{Studi Kasus dan Presentasi Proyek Akhir}

Praktikum terakhir semester ini sekaligus menjadi tempat studi kasus
dan presentasi proyek akhir. Setiap kelompok memilih satu program
dari bab-bab sebelumnya, lalu menjalankan seluruh alur kerja bab ini
terhadap program tersebut.

\begin{studikasusbox}
\textbf{Studi kasus: ``Program mahasiswa yang lambat''.} Sebuah
program pengolah larik buatan mahasiswa memerlukan waktu yang jauh
lebih lama daripada perkiraan pemiliknya. Program itu benar dan
keluarannya tepat, sehingga tidak ada kesalahan yang dapat ditemukan
dengan \eng{debugging} biasa. Tugas kelompok adalah menemukan
penyebabnya dengan perintah \code{U}, menunjukkannya melalui metrik,
memperbaikinya, dan membuktikan perbaikan itu dengan pengukuran.

\vspace{0.4em}
\textbf{Panduan singkat presentasi (7 menit + 3 menit tanya jawab):}
\begin{enumerate}
    \item Tampilkan potongan hasil \eng{disassembly} sebelum dan
          sesudah, dan tunjuk satu instruksi yang menjadi biang
          kelambatan (2 menit).
    \item Jelaskan teknik optimasi yang dipilih beserta syarat
          berlakunya (2 menit).
    \item Tampilkan tabel pengukuran sebelum dan sesudah; sebutkan
          metrik yang dipakai dan asumsi pengukurannya (2 menit).
    \item Nyatakan batas pekerjaan: apa yang \emph{tidak} berubah
          dan risiko apa yang masih tersisa (1 menit).
\end{enumerate}
Semua anggota kelompok harus mengambil bagian berbicara, dan setiap
klaim perbaikan wajib disertai angka.
\end{studikasusbox}

\begin{asesmenbox}
\textbf{Asesmen praktikum 15 (bobot 15\% dari komponen praktikum).}

\begin{tabular}{@{}p{6.0cm}cp{4.8cm}@{}}
    \toprule
    \textbf{Indikator} & \textbf{Bobot} & \textbf{Kriteria} \\
    \midrule
    Ketepatan hasil \eng{disassembly} & 25\% &
    Alamat, byte mesin, dan \eng{mnemonic} tercatat benar serta pola
    dikenali dengan tepat \\
    \addlinespace
    Kualitas identifikasi inefisiensi & 20\% &
    Minimal tiga inefisiensi ditemukan dan dikaitkan dengan metrik \\
    \addlinespace
    Kebenaran versi perbaikan & 25\% &
    Program berjalan, hasil \hex{1130} tetap, dan minimal tiga
    teknik optimasi diterapkan sesuai syaratnya \\
    \addlinespace
    Keabsahan pengukuran & 20\% &
    Masukan dan jumlah iterasi sama; pengukuran diulang dan nilai
    tengah dilaporkan \\
    \addlinespace
    Kerapian laporan dan jurnal & 10\% &
    Langkah kerja dapat diulang oleh pembaca lain \\
    \bottomrule
\end{tabular}
\end{asesmenbox}

\begin{asesmenbox}
\textbf{Rubrik presentasi proyek akhir (bobot 10\% dari komponen
praktikum).}

\begin{tabular}{@{}p{4.2cm}p{1.6cm}p{5.0cm}@{}}
    \toprule
    \textbf{Aspek} & \textbf{Bobot} & \textbf{Kriteria penilaian} \\
    \midrule
    Penguasaan hasil \eng{disassembly} & 25\% &
    Dapat menjelaskan makna byte dan pola tiap instruksi yang
    ditampilkan \\
    \addlinespace
    Argumentasi optimasi & 25\% &
    Setiap perubahan disertai alasan, syarat berlaku, dan
    kesadaran akan trade-off \\
    \addlinespace
    Bukti pengukuran & 25\% &
    Tabel sebelum/sesudah lengkap; keadilan pengukuran dijelaskan \\
    \addlinespace
    Penyajian dan tanya jawab & 25\% &
    Sistematis, waktu terpenuhi, seluruh anggota berbicara, jawaban
    tepat sasaran \\
    \bottomrule
\end{tabular}
\end{asesmenbox}

\section{Latihan}
\label{sec:latihan-bab16}

\begin{enumerate}
    \item \textbf{Konseptual.} Jelaskan mengapa \eng{disassembly}
          tidak dapat mengembalikan kode sumber asli, walaupun
          \eng{assembly} dan kode mesin bersifat satu-ke-satu.
          Sebutkan tiga informasi yang hilang selama perakitan dan
          jelaskan mengapa informasi itu tidak tersimpan di dalam
          berkas \file{.exe}.

    \item \textbf{Konseptual.} Byte \hex{48} dapat berarti instruksi
          \asm{DEC AX} atau karakter \code{'H'}. Jelaskan bagaimana
          keadaan ini menyulitkan \eng{disassembly}, lalu sebutkan
          dua cara praktis untuk menentukan batas antara daerah kode
          dan daerah data pada sebuah program.

    \item \textbf{Analisis hasil disassembly.} Perhatikan potongan
          berikut.

\begin{lstlisting}[language=MASM, numbers=none]
1D2F:0200 B90A00    MOV  CX,000A
1D2F:0203 31DB      XOR  BX,BX
1D2F:0205 0304      ADD  AX,[SI]
1D2F:0207 83C602    ADD  SI,+02
1D2F:020A E2F9      LOOP 0205
1D2F:020C 8906F401  MOV  [01F4],AX
\end{lstlisting}

          \begin{enumerate}
              \item Pola apa yang tampak pada potongan tersebut?
              \item Berapa kali instruksi \lstasm|ADD AX, [SI]|
                    dieksekusi? Jelaskan cara Anda memperolehnya.
              \item Berapa jarak relatif yang disimpan pada byte
                    \hex{E2 F9}, dan mengapa \eng{disassembler} dapat
                    menampilkan alamat \hex{0205}?
          \end{enumerate}

    \item \textbf{Analisis metrik.} Sebuah program memuat perulangan
          yang dijalankan 20 kali dengan isi
          \lstasm|ADD AX, [SI]|, \lstasm|ADD SI, 2|, dan
          \lstasm|LOOP|. Hitung perkiraan jumlah siklus seluruh
          perulangan dengan memakai Tabel~\ref{tab:siklus-8086-perkiraan}.
          Tuliskan rincian perhitungan Anda, dan sebutkan dua alasan
          mengapa angka hasil hitungan belum tentu sama dengan waktu
          sesungguhnya.

    \item \textbf{\eng{Strength reduction}.} Tuliskan hasil
          penggantian untuk potongan berikut, lalu sebutkan syarat
          yang harus dipenuhi agar penggantian itu sahih.

\begin{lstlisting}[language=MASM, numbers=none]
    MOV AX, 0
    ADD AX, nilai
    MOV BX, 16
    MUL BX
    CMP AX, 0
    JE  kosong
\end{lstlisting}

    \item \textbf{Optimasi dan trade-off.} Diberikan dua varian
          pengalian dengan empat: (a) \lstasm|MOV CL, 2| diikuti
          \lstasm|SHL AX, CL|, dan (b) dua kali \lstasm|SHL AX, 1|.
          Hitung ukuran byte dan perkiraan siklus keduanya. Varian
          mana yang Anda pilih bila program harus muat pada memori
          64 byte? Varian mana yang Anda pilih bila program berjalan
          di dalam perulangan seribu kali? Jelaskan alasan tiap
          jawaban.

    \item \textbf{Pengukuran.} Rancang pengukuran untuk
          membandingkan dua versi program yang mengerjakan pekerjaan
          sama. Sebutkan minimal empat syarat agar perbandingan itu
          adil, lalu tuliskan potongan kode yang mencatat waktu mulai
          dan waktu berhenti dengan \code{INT 1Ah} fungsi \hex{00}.
          Jelaskan mengapa percobaan pertama biasanya dibuang.

    \item \textbf{Penulisan program.} Tulislah program lengkap yang
          mengalikan setiap elemen larik \code{DW 3, 5, 7, 9, 11}
          dengan empat, lalu menjumlahkan seluruh hasil perkaliannya
          ke dalam \reg{DX}. Gunakan \asm{SHL}, bukan \asm{MUL}.
          Sertakan penghitung iterasi pada variabel tersendiri agar
          jumlah iterasi dapat diverifikasi, lalu periksa kebenaran
          hasilnya secara manual.
\end{enumerate}

\begin{refleksibox}
Jawab pertanyaan berikut pada jurnal praktikum Anda.
\begin{enumerate}
    \item Sebelum praktikum ini, apakah Anda pernah merasa yakin
          bahwa sebuah perubahan pasti mempercepat program? Setelah
          mengukur, apakah dugaan Anda terbukti? Bagian mana yang
          paling berbeda dari dugaan?
    \item Teknik optimasi mana yang paling sulit Anda terapkan
          dengan benar, dan kesalahan seperti apa yang pernah Anda
          buat saat menerapkannya?
    \item Menurut Anda, pada situasi seperti apa sebuah program
          sebaiknya dibiarkan dalam bentuk naif saja? Kaitkan
          jawaban Anda dengan nilai keterbacaan dan kemudahan
          pemeliharaan.
\end{enumerate}
\end{refleksibox}

\section{Rangkuman}
\label{sec:rangkuman-bab16}

Pada bab ini telah dibahas:

\begin{itemize}
    \item \eng{Disassembly} adalah pekerjaan mengubah kode mesin
          kembali menjadi \eng{mnemonic}. Pekerjaan ini searah dan
          tidak dapat membalikkan \eng{assembly} sepenuhnya, karena
          nama simbolik, komentar, dan batas data hilang saat
          perakitan.

    \item Perintah \code{U} pada DEBUG menampilkan alamat, byte
          mesin, dan \eng{mnemonic} sekaligus. Perkakas lain yang
          sejenis adalah \eng{disassembler} Emu8086, GUI Turbo
          Assembler, \code{ndisasm}, dan \code{objdump}.

    \item Keterbatasan utama \eng{disassembly} adalah sulitnya
          membedakan kode dan data, munculnya alamat relatif dan
          \eng{displacement} menggantikan nama variabel, serta
          hilangnya semua penanda yang berguna bagi manusia.

    \item Pola hasil \eng{disassembly} dapat dikenali: pembandingan
          dan percabangan, perulangan, \eng{prologue} dan
          \eng{epilogue} prosedur, inisialisasi segmen, serta
          pemanggilan interupsi.

    \item Empat metrik optimasi adalah jumlah instruksi yang
          dieksekusi, perkiraan jumlah siklus, ukuran kode dalam
          byte, dan jumlah akses memori. Tanpa pengukuran, optimasi
          hanyalah dugaan yang belum terbukti.

    \item Teknik optimasi level instruksi meliputi: memilih register
          daripada memori, mengganti \asm{MUL}/\asm{DIV} dengan
          pergeseran, memakai \asm{LEA} untuk aritmatika yang tidak
          boleh mengubah \eng{flag}, mengurangi lompatan,
          \eng{loop unrolling}, dan \lstasm|XOR reg, reg| sebagai
          pengganti \lstasm|MOV reg, 0|.

    \item \eng{Strength reduction} menggantikan operasi mahal dengan
          operasi murah, tetapi hanya sahih bila syaratnya dipenuhi,
          misalnya pengali harus pangkat dua dan nilai harus tak
          bertanda.

    \item Optimasi ukuran dan optimasi kecepatan sering bertentangan.
          Pilihan ditentukan oleh batasan proyek, bukan oleh selera.

    \item Pengukuran yang adil menuntut masukan sama, jumlah iterasi
          sama, batas pengukuran yang sama, dan pengulangan beberapa
          kali.

    \item Batas etika optimasi adalah kebenaran lebih dahulu,
          larangan optimasi dini, serta kewajiban menjaga
          keterbacaan dan dokumentasi.
\end{itemize}

\begin{checklistbox}
Tandai setiap butir yang sudah Anda kuasai.

\begin{itemize}
    \item[\checkbox] Saya dapat menjelaskan perbedaan
          \eng{assembly} dan \eng{disassembly} beserta keterbatasan
          \eng{disassembly}.
    \item[\checkbox] Saya dapat membongkar program dengan perintah
          \code{U} dan membaca keluaran alamat, byte, serta
          \eng{mnemonic}.
    \item[\checkbox] Saya dapat mengenali pola perulangan,
          percabangan, \eng{prologue} prosedur, dan pemanggilan
          interupsi pada hasil \eng{disassembly}.
    \item[\checkbox] Saya dapat menghitung jumlah instruksi yang
          dieksekusi, ukuran byte, perkiraan siklus, dan jumlah akses
          memori.
    \item[\checkbox] Saya dapat menerapkan minimal empat teknik
          optimasi level instruksi beserta syarat berlakunya.
    \item[\checkbox] Saya dapat menjelaskan trade-off antara optimasi
          ukuran dan optimasi kecepatan.
    \item[\checkbox] Saya dapat merancang pengukuran yang adil dan
          mencatat waktu dengan \code{INT 1Ah}.
    \item[\checkbox] Saya dapat menjelaskan alasan mengapa optimasi
          dini harus dihindari.
\end{itemize}

\vspace{0.4em}
Bila masih ada butir yang belum bertanda, ulangi bagian yang
bersesuaian sebelum mengerjakan evaluasi.
\end{checklistbox}

\begin{infobox}
\textbf{Rujukan bab.} Materi pada bab ini diperkuat oleh
\cite{wikibooksdisasm}, \cite{bartlett2023}, dan \cite{nasm2024}.
Untuk pembahasan yang lebih mendalam mengenai keterbatasan
\eng{disassembly} dan cara menafsirkan keluaran perkakas pembongkar,
lihat \cite{hyde2010}.
\end{infobox}

\section{Evaluasi}
\label{sec:evaluasi-bab16}

\subsection{Soal Pilihan Ganda}

\begin{enumerate}
    \item Perintah pada MS-DEBUG yang dipakai untuk membongkar kode
          mesin menjadi \eng{mnemonic} adalah \dots
          \begin{enumerate}
              \item \code{T} \item \code{U}
              \item \code{D} \item \code{G}
          \end{enumerate}

    \item Nama variabel \code{larik} tidak muncul pada hasil
          \eng{disassembly}, melainkan berubah menjadi
          \lstasm|[SI+01F4]|. Penyebabnya adalah \dots
          \begin{enumerate}
              \item nama simbolik tidak disimpan di dalam berkas biner
              \item DEBUG selalu menghapus nama simbolik
              \item linker mengubah nama simbolik menjadi alamat
                    relatif
              \item disassembler hanya mampu membaca register
          \end{enumerate}

    \item Instruksi tunggal yang tepat menggantikan
          \lstasm|MUL BX| untuk mengalikan \reg{AX} (tak bertanda)
          dengan delapan adalah \dots
          \begin{enumerate}
              \item \lstasm|SHR AX, 3| \item \lstasm|SHL AX, 3|
              \item \lstasm|SAR AX, 3| \item \lstasm|DIV AX, 8|
          \end{enumerate}

    \item Ciri paling dapat diandalkan bahwa suatu blok pada hasil
          \eng{disassembly} adalah perulangan adalah \dots
          \begin{enumerate}
              \item adanya instruksi \asm{INT 21h}
              \item lompatan mundur ke alamat yang lebih kecil
              \item adanya \lstasm|MOV AX, @DATA|
              \item lompatan maju yang jauh
          \end{enumerate}

    \item Keunggulan \lstasm|XOR AX, AX| dibandingkan
          \lstasm|MOV AX, 0| untuk mengosongkan register adalah \dots
          \begin{enumerate}
              \item hasilnya selalu nol pada semua kondisi
              \item lebih pendek satu byte dan lebih cepat satu
                    siklus
              \item aman dipakai di antara \asm{CMP} dan lompatan
                    bersyarat
              \item sekaligus mengosongkan \reg{DX}
          \end{enumerate}
\end{enumerate}

\subsection{Soal Uraian}

\begin{enumerate}
    \item Jelaskan tiga keterbatasan \eng{disassembly}. Untuk
          masing-masing keterbatasan, berikan satu contoh potongan
          byte atau instruksi yang memperlihatkan mengapa
          \eng{disassembly} sulit ditafsirkan.

    \item Perhatikan potongan hasil \eng{disassembly} berikut.

\begin{lstlisting}[language=MASM, numbers=none]
1D2F:0100 B90A00    MOV  CX,000A
1D2F:0103 B80000    MOV  AX,0000
1D2F:0106 0304      ADD  AX,[SI]
1D2F:0108 83C602    ADD  SI,+02
1D2F:010B E2F9      LOOP 0106
1D2F:010D D1E0      SHL  AX,1
1D2F:010F D1E0      SHL  AX,1
1D2F:0111 D1E0      SHL  AX,1
1D2F:0113 8906F401  MOV  [01F4],AX
\end{lstlisting}

          \begin{enumerate}
              \item Sebutkan pola yang tampak dan tunjukkan buktinya.
              \item Hitung jumlah instruksi yang dieksekusi untuk
                    seluruh program tersebut.
              \item Hitung perkiraan jumlah siklusnya, lalu
                    sebutkan satu perubahan pada potongan itu yang
                    masih dapat menghemat siklus.
          \end{enumerate}

    \item Diberikan program naif berikut: akumulator di memori,
          pengali 16 dimuat ke \reg{BX} setiap iterasi, perkalian
          memakai \asm{MUL}, dan penghitung perulangan memakai
          \asm{DEC} diikuti \asm{JNZ}. Tuliskan versi teroptimasi
          lengkap untuk program tersebut, sebutkan nama setiap
          teknik yang Anda pakai, dan jelaskan syarat berlaku
          masing-masing. Sebutkan pula metrik yang Anda pakai untuk
          membuktikan perbaikannya.

    \item Seorang mahasiswa mengganti sebuah \asm{MUL} dengan
          pergeseran, lalu melaporkan bahwa programnya menjadi dua
          kali lebih cepat. Rancang langkah-langkah pengukuran yang
          harus ia lakukan agar klaim tersebut dapat dipercaya,
          kemudian jelaskan alasan mengapa optimasi dini justru
          lebih sering merugikan daripada menguntungkan.
\end{enumerate}

\subsection{Kisi-kisi Penilaian}

\begin{table}[H]
    \centering
    \caption{Kisi-kisi evaluasi Bab~\ref{chap:disassembly}.}
    \label{tab:kisi-bab16}
    \small
    \begin{tabular}{@{}p{1.5cm}p{5.0cm}p{2.6cm}c@{}}
        \toprule
        \textbf{Sub-CPMK} & \textbf{Indikator yang diukur} &
        \textbf{Bentuk soal} & \textbf{Bobot} \\
        \midrule
        4.3 & Membaca hasil \eng{disassembly} dan mengenali polanya &
        Pilihan ganda 1, 2, 4; Uraian 1, 2 & 30\% \\
        \addlinespace
        4.3 & Menjelaskan keterbatasan \eng{disassembly} &
        Pilihan ganda 2; Uraian 1 & 15\% \\
        \addlinespace
        4.3 & Menghitung metrik optimasi &
        Uraian 2 bagian b dan c & 20\% \\
        \addlinespace
        4.3 & Menerapkan teknik optimasi level instruksi &
        Pilihan ganda 3, 5; Uraian 3 & 25\% \\
        \addlinespace
        4.3 & Merancang pengukuran dan memahami batas optimasi &
        Uraian 4 & 10\% \\
        \bottomrule
    \end{tabular}
\end{table}

\begin{checklistbox}
\textbf{Checklist Pencapaian Kompetensi Bab~\ref{chap:disassembly}}

Seorang mahasiswa dinyatakan kompeten pada Sub-CPMK 4.3 apabila:
\begin{itemize}
    \item[\checked] dapat membongkar program dengan perintah
          \code{U} serta membaca alamat, byte mesin, dan
          \eng{mnemonic} dengan benar;
    \item[\checked] dapat menyebutkan keterbatasan \eng{disassembly}
          dan menjelaskan mengapa kode serta data sulit dibedakan;
    \item[\checked] dapat mengenali minimal lima pola hasil
          \eng{disassembly} beserta cirinya;
    \item[\checked] dapat menghitung empat metrik optimasi dan
          menjelaskan bahwa tanpa pengukuran optimasi hanyalah
          dugaan;
    \item[\checked] dapat menerapkan minimal tiga teknik optimasi
          level instruksi beserta syarat berlakunya;
    \item[\checked] dapat menjelaskan trade-off ukuran dan kecepatan
          serta memilih dengan alasan;
    \item[\checked] berhasil membuktikan perbaikan program melalui
          pengukuran yang adil.
\end{itemize}
\end{checklistbox}


\end{document}
